
\chapter{Appendix: Sustianable Development Goals}
The analysis and separation of electromyographic (EMG) signals developed
throughout this Bachelor's Thesis not only aims to contribute to advances in the
field of biomedical signal processing, but also naturally aligns with several
Sustainable Development Goals, also known as
\href{https://sdgs.un.org/goals}{SDGs}. The application of techniques such as
preprocessing, artifact removal and separation and classification methods
contributes to the development of technological solutions that improve health,
promote innovation and support equal access to biomedical tools.

Therefore, the main SDGs addressed by this Bachelor's Thesis and their role in
the project are presented below:

\begin{itemize}
    \item Goal 3: Good Health and Well-being
\end{itemize}

Improving the quality of the acquired signals enables more accurate clinical
assessments, supports the early detection of muscular dysfunctions and optimizes
rehabilitation processes. The development of more robust and reliable algorithms
has a direct impact on the quality of healthcare by supporting decision-making
among healthcare professionals.

\begin{itemize}
    \item Goal 9: Industry, Innovation and Infrastructure
\end{itemize}

The project promotes technological innovation by applying engineering tools, such as digital signal processing and machine learning, to a
biomedical context. The study and comparison of separation algorithms reinforce
the development of new tools that can be integrated into medical devices prostheses or monitoring systems. Therefore, this work
contributes to building scientific and technological infrastructures to improve health.

\begin{itemize}
    \item Goal 10: Reduced Inequalities
\end{itemize}

By advancing techniques that make it possible to obtain accurate muscular
information without relying on high-cost equipment, this project supports the
development of more accessible solutions. It favors the extension of EMG technologies to clinical and research environments with fewer available
resources. The standardization and validation of separation methods allow these tools to be reproduced more reliably, facilitating their use in a wider variety of centers. In this way, the project promotes more equitable access to prosthetic technologies, contributing to the reduction of inequalities in the
healthcare field.